% !PRINTED 711
\weberpage{711}
%<<CONT
ist, und danach sind auch die Functionen $f$ in $f_1$ und $f_2$ zu
unterscheiden, so dass

\weq{11}{f_1(r^{-s}) = f_1(r^s), \quad f_2(r^{-s}) = -\,f_2(r^s)}
wird. Ebenso unterscheiden wir die Ausdrücke (8), (9) als $X_1$
und $X_2$, und es ergiebt sich, wenn man beide Ausdrücke addirt,
nach (11):
\weq{12}{\begin{aligned}
X_1 &= -\frac{1}{2m}\sum_{1,\,m-1}^{s} f_1(r^s)\,\log\left(\sin\frac{s\pi}{m}\right)^{2}, \\
X_2 &= -\frac{i\pi}{m^2}\sum_{1,\,m-1}^{s} s\,f_2(r^s).
\end{aligned}}

\begin{center}
\S.\,195.\\[4pt]
\gesperrt{Ueber die Classenzahl in dem in $\Omega_m$ enthaltenen
reellen Körper.}
\end{center}

Unsere allgemeinen Principien können auch angewandt werden, um die
Classenzahlen in den Theilern des Körpers $\Omega_m$ zu bestimmen.

Insbesondere sind die Formeln (6), (7), \S.\,192 anwendbar, wenn
wir die Grade der Primideale in einem solchen Theiler kennen. Die
Frage nach dem Verhältniss der Classenzahlen in den Theilern eines
Körpers $\Omega_m$ zu der Classenzahl in $\Omega_m$ selbst ist von
grossem Interesse und soll hier für den Fall behandelt werden, dass
es sich um den in $\Omega_m$ enthaltenen reellen Körper $H_m$
handelt, den wir ja schon im \S.\,176 f. untersucht haben\footnote{Vgl.
\gesperrt{Kummer}, \glqq Bestimmung der Anzahl u.\,s.\,f.\grqq{}.
\gesperrt{Crelle's Journ.}, Bd.\,40 (1849).}. Es hat sich dort ergeben, dass
$q$ im Körper $H_m$ die $\tfrac12\mu^{\text{te}}$ Potenz einer
Primzahl ersten Grades ist, dass die anderen Primzahlen $p$ in $e_1$
Primfactoren $f_1^{\text{ten}}$ Grades zerfallen, so dass

\weq{1}{f_1 e_1 = \tfrac{1}{2}\mu,}
und dass zwei Arten von Primzahlen $p_1$, $p_2$ zu unterscheiden
sind, so dass bei den Primzahlen erster Art $f_1 = \tfrac12 f$,
$e_1 = e$, bei denen der zweiten Art $f_1 = f$, $e_1 = \tfrac12 e$
ist.

Für den Körper $H_m$ erhalten wir demnach aus \S.\,192, (6), (7)
die Classenzahl $h_1$ nach der Formel:
%CONT>>
